\documentclass[a4paper,12pt]{article}

% Paquetes necesarios
\usepackage[utf8]{inputenc}   % Codificación UTF-8
\usepackage{fontspec} % Paquete para manejar fuentes específicas
\usepackage[spanish]{babel}
\usepackage[style=apa, backend=biber]{biblatex}
\usepackage{geometry}         % Configuración de márgenes
\usepackage{graphicx}
\usepackage{setspace}         % Configuración de interlineado
\usepackage{fancyhdr}         % Configuración de encabezados y pie de página
\usepackage{biblatex}
\usepackage{hyperref}

\setmainfont{Arial}   % Configura Arial como la fuente principal
\geometry{
    top=2.5cm,
    bottom=2.5cm,
    left=4cm,
    right=2.5cm
}
\pagestyle{fancy}
% Configuración del interlineado y espacio entre párrafos
\onehalfspacing            % Interlineado de 1.5
% Configuración del pie de página
\fancyhf{}
\fancyfoot[R]{\thepage}    % Número de página en la esquina inferior derecha
\fancypagestyle{plain}{
    \fancyhf{}
    \fancyfoot[R]{\thepage}
}

\addbibresource{bibliography.bib} % Archivo .bib con las referencias


\begin{document}

% Incluir la portada
% Archivo: portada.tex
\begin{titlepage}
    \begin{center}
        % Logo
        \includegraphics[width=0.8\textwidth]{images/logo_soe.png} \\
        \vspace{1cm}
        % Título de la institución

        {\fontsize{18pt}{10pt}\selectfont\textbf{FACULTAD DE INGENIERÍA EN CIENCIAS DE LA COMPUTACIÓN Y TELECOMUNICACIONES}} \\
        \vspace{2cm}
        {\fontsize{18pt}{5pt}\selectfont\textbf{UNIDAD DE POSTGRADO SCHOOL OF ENGINEERING}} \\
        \vspace{0.5cm}
        
        % Nombre de la maestría
        \vspace{0.5cm}
        {\fontsize{20pt}{5pt}\selectfont\textbf{MONOGRAFIA}} \\
        \vspace{1cm}
        
        % Título del trabajo
        {\fontsize{16pt}{5pt}\selectfont\textbf{Anáslis de programación multiplataforma para abaratar costos en el desarrollo de aplicaciones móbiles}} \\
        \vspace{1cm}
        
        \vspace{2cm}
        {\fontsize{14pt}{22pt}\selectfont\textbf{Ing. Joel Gabriel Torrejón Méndez}} \\
        
        % Ciudad y fecha
        \vfill
        \textbf{Santa Cruz, Bolivia} \\
    \end{center}
\end{titlepage}
     % Reemplaza con % Archivo: portada.tex
\begin{titlepage}
    \begin{center}
        {\fontsize{15pt}{10pt}\selectfont\textbf{UNIVERSIDAD AUTÓNOMA "GABRIEL RENE MORENO"}} \\
        {\fontsize{15pt}{10pt}\selectfont\textbf{FACULTAD DE INGENIERÍA EN CIENCIAS DE LA COMPUTACIÓN Y TELECOMUNICACIONES}} \\
        \vspace{1cm}
        % Logo
        \includegraphics[width=0.8\textwidth]{images/logo_soe.png} \\
        
        % Título de la institución
        {\fontsize{15pt}{10pt}\selectfont\textbf{MAESTRÍA EN CIENCIA DE DATOS E INTELIGENCIA ARTIFICIAL}} \\

        % Título del trabajo
        \vspace{1cm}
        {\fontsize{15pt}{5pt}\selectfont\textbf{Implementar un Chatbot Semántico Multiperfil para Acceso Inteligente a Documentación Técnica y de Negocio en Plataformas Confluence del Sector Aeronáutico en Airnguru S.A. para la Gestión 2025}} \\
        \vspace{1cm}
        
        \vspace{2cm}
        \begin{flushright}
        {\fontsize{14pt}{22pt}\selectfont\textbf{Autor}} \\
        {\fontsize{14pt}{22pt}\selectfont{Ing. Darlyn Bravo Peña}} \\
        \end{flushright}
        
        
        
        % Ciudad y fecha
        \vfill
        \textbf{Santa Cruz, Bolivia} \\
    \end{center}
    \end{titlepage}
     si prefieres que se trate como una sección independiente

% Table of contents
\newpage
\tableofcontents
% \listoffigures

\section{Introducción}

La programación híbrida en aplicaciones móviles es una estrategia ideal para desarrollar
soluciones multiplataforma con un presupuesto reducido. Este enfoque permite crear una
única base de código utilizando tecnologías web como HTML, CSS y JavaScript, que luego
se empaqueta en un contenedor nativo para ejecutarse en dispositivos Android e iOS.
Frameworks como Ionic, React Native o Flutter optimizan el proceso, permitiendo un
desarrollo ágil y reduciendo significativamente los costos y tiempos asociados con la
creación de apps separadas para cada plataforma.\\

Además de abaratar costos iniciales, las aplicaciones híbridas simplifican el mantenimiento
y las actualizaciones, ya que los cambios se realizan sobre un solo código y se reflejan en
todas las plataformas. Aunque pueden tener limitaciones en cuanto al rendimiento para
aplicaciones complejas, son ideales para startups o empresas que necesitan una app
funcional, atractiva y accesible en ambos sistemas operativos sin invertir en desarrolladores
especializados para cada ecosistema.\\

El tema está ubicado en el area 
Ciencias Generales de la Computación Aplicadas y la 
línea Ingeniería de software con un eje de 
Aspectos económicos y de negocio del proceso de desarrollo del software porque los costos
de desarrollo de software son un factor crítico en la toma de decisiones de las empresas.

\section{Antecedentes}

Con el aumento significativo de usuarios de teléfonos móviles en la última
década, casi todas las empresas necesitan desarrollar una aplicación móvil.
Sin embargo, uno de los principales factores que debe tener en cuenta al
desarrollar una aplicación móvil es elegir entre el desarrollo de aplicaciones
móviles nativas o híbridas. \parencite{turing-hibrid-development} \\

Si bien ambas soluciones tienen pros y contras, muchas empresas optan por
las aplicaciones híbridas debido a sus beneficios en términos de presupuesto,
tiempo de comercialización y otros factores. Según un informe de Forbes \parencite{forbes-hibrid-development},
37 de las 50 principales aplicaciones minoristas en Estados Unidos son híbridas.
Además, plataformas populares como Twitter, Instagram, Gmail, Uber, etc.,
utilizan aplicaciones híbridas.\\

En respuesta a estas limitaciones, surgió la programación híbrida como una
alternativa más eficiente y accesible. Este enfoque permite desarrollar
aplicaciones utilizando tecnologías web estándares como HTML, CSS y JavaScript,
combinadas con frameworks que facilitan su integración con funcionalidades
nativas de los dispositivos. Frameworks como Ionic, React Native y Flutter
han demostrado su capacidad para reducir significativamente los costos
y los tiempos de desarrollo, manteniendo una experiencia de usuario satisfactoria.
Este proyecto se centra en analizar cómo la programación híbrida puede ser
una solución viable para optimizar recursos en el desarrollo de aplicaciones
móviles, especialmente en proyectos donde el presupuesto es un factor crítico.

\newpage
\section{Planteamiento del problema}

\subsection{Presentación de las problemáticas}
El desarrollo de aplicaciones móbiles es parte fundamental para las
empresas que buscan expandir su mercado y llegar a un público más amplio. Esto
se debe a que la mayoría de las personas cuentan con un dispositivo móvil y\
acceden a internet a través de él. Sin embargo, el desarrollo de aplicaciones
enfrenta problemáticas como:
\begin{itemize}
    \item \textbf{Mantenimiento de aplicaciones:} Las aplicaciones móviles
        requieren de mantenimiento constante para corregir errores y agregar
        nuevas funcionalidades.
    \item \textbf{Precios elevados:} El desarrollo de aplicaciones móviles
        puede ser costoso, especialmente para pequeñas y medianas empresas.
    \item \textbf{Dificultad para encontrar desarrolladores:} Encontrar
        desarrolladores con experiencia en el desarrollo de aplicaciones
        móviles puede ser complicado.
    \item \textbf{Dificultad para mantenerse actualizado:} El desarrollo de
        aplicaciones móviles requiere de estar al tanto de las últimas
        tecnologías y tendencias.
\end{itemize}

\subsection{Situación problemática}
El desarrollo de aplicaciones móviles es un proceso complejo que requiere de
conocimientos especializados y experiencia. Además, el desarrollo de
aplicaciones móviles puede ser costoso y requiere de mantenimiento constante
para corregir errores y agregar nuevas funcionalidades.\\

Estas limitantes agregan una barrera de entrada en el mercado de aplicaciones
móviles, especialmente para pequeñas y medianas empresas que no cuentan con
los recursos necesarios para desarrollar y mantener una aplicación móvil.

\subsection{Pregunta de investigación}
¿ Como puede el uso de tecnologías de desarrollo multiplataforma ayudar
a simplificar el desarrollo de aplicaciones móviles y abaratar costos?


\newpage
\section{Justificación}

\subsection{Justificación teórica}
El desarrollo de aplicaciones móviles enfrenta desafíos económicos
significativos debido a la necesidad de adaptar las soluciones a diferentes
plataformas, como Android e iOS. El enfoque tradicional de desarrollo nativo
implica la creación de aplicaciones separadas para cada sistema operativo,
utilizando lenguajes y herramientas específicas. Esto no solo incrementa los
costos relacionados con la contratación de desarrolladores especializados, sino
también los tiempos de desarrollo y mantenimiento.\\

En este contexto, el análisis de la programación multiplataforma emerge como
una alternativa teórica y práctica que busca mitigar estos desafíos
económicos.\\


\subsection{Justificación práctica}
Desde una perspectiva práctica, la programación multiplataforma ofrece
una solución directa a los problemas asociados con el desarrollo nativo
al reducir costos, tiempos y complejidad operativa. Empresas y
desarrolladores han adoptado frameworks multiplataforma debido a su
capacidad para generar aplicaciones funcionales y visualmente
atractivas a partir de un único código base.\\

Este enfoque elimina la necesidad de mantener equipos especializados
en múltiples lenguajes y herramientas, disminuyendo significativamente
el gasto en recursos humanos y tecnológicos.

\newpage
\section{Objetivos}

\subsection{Objetivo general}
Analizar el impacto de la programación multiplataforma como estrategia
para reducir costos en el desarrollo de aplicaciones móviles, evaluando
su viabilidad técnica y económica, así como su capacidad para satisfacer
las necesidades de funcionalidad y experiencia de usuario en diferentes
plataformas.

\subsection{Objetivos específicos}
\begin{itemize}
    \item \textbf{Investigar} las características y principios fundamentales de la programación multiplataforma, así como los principales frameworks utilizados en el desarrollo de aplicaciones móviles.
    \item \textbf{Identificar} los factores económicos, técnicos y operativos que influyen en los costos del desarrollo nativo y multiplataforma.
    \item \textbf{Comparar} el rendimiento, la funcionalidad y la experiencia de usuario entre aplicaciones desarrolladas de forma nativa y multiplataforma.
    \item \textbf{Evaluar} la eficiencia del desarrollo multiplataforma en términos de reducción de tiempos, costos de desarrollo y mantenimiento.
    \item \textbf{Proponer} recomendaciones prácticas para la adopción de programación multiplataforma en proyectos móviles, considerando las necesidades específicas de empresas con recursos limitados.
\end{itemize}
\newpage
\section{Alcance}

\subsection{Delimitación o Alcance especial}
El presente estudio se enfocará en analizar el impacto de la programación
multiplataforma en equipos de desarrollo de aplicaciones móviles ubicados
en la ciudad de Santa Cruz. Se tomarán como referencia empresas tecnológicas,
startups y desarrolladores independientes que operan en esta región,
considerando sus necesidades específicas, recursos disponibles y prácticas
actuales de desarrollo.\\

Este enfoque permitirá obtener datos relevantes
y contextuales que reflejen las condiciones económicas y técnicas propias
del sector en Santa Cruz, facilitando la formulación de estrategias
adaptadas a su realidad local.

\subsection{Alcance temporal}
El análisis abarcará el período comprendido desde el año 2015 hasta la
actualidad. Este intervalo se selecciona debido a que marca el auge y
evolución de frameworks de programación multiplataforma, los cuales
han transformado significativamente el desarrollo de aplicaciones móviles.\\

Este enfoque temporal permitirá evaluar cómo ha cambiado la adopción de
estas tecnologías y su impacto en los costos y la eficiencia del
desarrollo en el contexto de los equipos de Santa Cruz.



\newpage
\printbibliography

\end{document}
